\documentclass[11pt]{article}
\usepackage[a4paper,margin=26mm]{geometry}
\usepackage{amsmath,amssymb,amsthm,mathtools,booktabs,longtable,hyperref}
\usepackage[T1]{fontenc}
\hypersetup{hidelinks}
\newtheorem{theorem}{Theorem}
\newtheorem{lemma}[theorem]{Lemma}
\newtheorem{proposition}[theorem]{Proposition}
\newtheorem{corollary}[theorem]{Corollary}
\newtheorem{remark}[theorem]{Remark}
\newcommand{\TV}{\operatorname{TV}}
\newcommand{\KL}{\operatorname{KL}}
\newcommand{\sig}{\sigma}
\newcommand{\eps}{\varepsilon}
\newcommand{\calE}{\mathcal E}
\newcommand{\calU}{\mathcal U}
\newcommand{\ind}{\mathbf 1}
\DeclareMathOperator{\logit}{logit}
\title{Active local acquisition of a rare correlated failure certificate\\\large A costed construction, physical preparation, preventive choice, and an amplitude--exposure obstruction}
\author{Astra research construction: independent intervention endpoint}
\date{5 November 2026}
\begin{document}
\maketitle
\begin{abstract}
We give an explicit experimental access model under which local experiments acquire an unknown correlated reliability mechanism and certify extremely rare system-level failures without requiring an observed system failure. Calibrated fields balance local binary conditional probabilities; finite-time observations and martingale concentration estimate local energy differences. Priced state pins then prepare a fast escape direction at each operating field, and measured dwell times recover the operating kinetic rates using local detailed balance. No invariance of the attempt clock across bias fields is assumed. The resulting rational model and confidence envelope compose with the previously accepted correlated-reliability theorem to yield relative stationary risk and a simultaneous relative first-failure CDF. A finite, physically generated preparation law replaces an exact equilibrium-safe initialization oracle. A finite menu of preventive fields can be compared with a multiplicative guarantee. Acquisition time, control amplitude, slew, state preparation, observation, and model enumeration are all charged. A controlled-trajectory lower bound shows that an amplitude cap below the compensating energy scale restores an exponential exposure cost, even with free parameter-independent resets. This is a proved composition under declared intervention and closure assumptions, not a claim to have validated those assumptions from data or established historical firstness.
\end{abstract}

\section{What is completed, and the inherited interface}
The endpoint is an acquisition-to-decision operation. It starts with a complete interaction graph and a physical local-control interface, but no supplied factor values, rare transition rates, stationary samples, or exact safe-equilibrium start. It ends with a finite confidence certificate and a preventive choice from a declared finite menu. The experiment does not count global catastrophes. Its physical powers, especially energy compensation and state pinning, are substantial and explicitly accounted for.

We use the accepted correlated-reliability theorem exactly as supplied, without re-proving it \cite{accepted,prior}. Its input is a positive rational factor model on $N$ bits, a visible network with $m$ edge bits, a fixed conditional-surprisal contrast $c$, certified Dobrushin contraction, and a reversible rational generator with a positive heat-bath backbone. It computes relative stationary disconnection probability and a simultaneous relative first-failure CDF from the model's own equilibrium distribution conditioned on safety, with polynomial work at fixed $c$ and no inverse global-rarity factor. We also use its accepted robustness interface: if total factor log error is $E$ and every nonzero transition-rate log error is at most $b$, on common support, then
\begin{equation}
 e^{-(6E+b)}\widehat F_u(t)\le F_u(t)\le e^{6E+b}\widehat F_u(t),\qquad t\ge0,
 \label{eq:inherited}
\end{equation}
where each law has its own equilibrium-safe start. Stationary probabilities of any event differ by at most $e^{\pm2E}$. These are accepted inputs to the present construction, not new claims about operator comparison.

The physical preparation operation in Section~\ref{sec:prep} is a separate addition: it produces a pointwise relative approximation to the true equilibrium-safe law using finite elapsed time. Absolute total variation alone would not give an all-time relative CDF bound near zero.

\section{Unknown mechanism and exact acquisition assumptions}
\label{sec:access}
Let $\Gamma=([N],\calE)$ be a known undirected interaction graph, with $M=|\calE|$ and maximum degree $d$. A separate known connected visible network has $m\le N$ edge bits. A visible bit equal to one means failed. The event $B$ is visible-network disconnection and $A=B^c$. Extra variables are declared and observed; no omitted common cause is silently absorbed into the model.

The unknown equilibrium energy has the pairwise form
\begin{equation}
 W_\theta(x)=\exp\left\{\sum_i h_ix_i+\sum_{\{i,j\}\in\calE}J_{ij}x_ix_j\right\},\qquad
 \ell_i(s)=h_i+\sum_{j\in\partial i}J_{ij}s_j,\qquad |\ell_i(s)|\le B_0.
 \label{eq:model}
\end{equation}
The known number $B_0$ measures an energy in units where inverse temperature is one. In particular $|h_i|\le B_0$ and $|J_{ij}|\le2B_0$. A finite menu $\calU$ of $K$ calibrated, reproducible, constant one-body fields is supplied. Under $u\in\calU$ the invariant law is proportional to $W_\theta(x)e^{u\cdot x}$. Fields may have certified readback errors, treated below.

The true operating generator has only single-coordinate flips:
\begin{equation}
 Q^u(x,x^{(i)})=r_i^u(x_{\partial i})\,
 \sig\bigl((1-2x_i)(\ell_i(x_{\partial i})+u_i)\bigr),\qquad
 \sig(a)=\frac1{1+e^{-a}},
 \label{eq:generator}
\end{equation}
with unknown local clocks $r_i^u(s)\in[r_*,r^*]$, where $0<r_*\le r^*<\infty$ are known. The clocks can depend arbitrarily on the neighbor context and operating field. Formula~\eqref{eq:generator} is equivalent to positive two-state local detailed balance with the stated invariant conditional distribution. It is a model assumption, not a consequence of observing an equilibrium histogram.

The apparatus has the following declared properties.
\begin{enumerate}
\item \textbf{Faithful local cells and clamps.} For a requested $i,s$, an isolated test cell or safely isolated copy can hold the neighbors of $i$ at $s$. Once held, the focal variable is a time-homogeneous two-state Markov chain with the same local energy difference as the operating system. For each requested menu setting $u$, the apparatus reproduces its corresponding local physical environment and focal field $u_i$, so its two rates are the actual operating rates in~\eqref{eq:generator}. This is required even when $r_i^u(s)$ depends on the full menu setting. There is no memory of the clamp once the requested condition is in place. Cell construction, isolation, setting the neighbor pattern, and settling to the declared Markov regime have stated finite costs, denoted collectively $C_{\rm cell}$ and $C_{\rm context}$. No equilibrium reset is supplied.
\item \textbf{Calibrated additive bias.} Under a held local bias $z$, the conditional success probability is $\sig(\ell_i(s)+z)$. The sum of the two biased rates can be an arbitrary unknown $a_i(s,z)\ge r_{\rm cal}>0$. No equality between this sum and an operating clock is assumed. The fixed value actually realized by a command is known within $\xi\le\eta/16$ and is reproducible during a hold. The permitted amplitude contains $[-B_0-\xi,B_0+\xi]$ and all operating fields. Commands have a specified calibration cost $C_{\rm field}(\xi)$ and a certified available slew schedule. When constant speed $v$ is available, travel $V$ takes $V/v$ elapsed time; otherwise the supplied schedule cost $C_{\rm slew}(V)$ is charged. A maximum speed alone gives only a lower bound on this elapsed cost. Calibration errors and drift beyond the stated intervals invalidate the theorem's premise.
\item \textbf{Priced state preparation.} At a held context and operating field, the focal bit can be pinned to a specified state and released. Each operation has a known elapsed cost $C_{\rm pin}(i,s,u)$, including preparation and settling. Actuator work is recorded as a separate resource and is not converted into time without an explicit apparatus model. The post-release first-exit clock follows the operating generator from that known state. This is active preparation of a possibly rare local state, not a free sample from the operating law. A physical implementation whose preparation cost is exponential pays that exponential cost.
\item \textbf{Observations and timing.} A focal bit can be read exactly at a stated per-read cost. First-exit events after release can be detected, with timestamp error at most $\Delta$ per endpoint, where $\Delta\le\rho/(64r^*)$. The event detector need only detect the first exit; subsequent motion is irrelevant. These are finite-resolution measurements. A whole operating configuration can be read for the final safety test at the completion of a known deterministic inspection duration. The reading refers to the state at that completion instant and is noninvasive; the mission clock begins at that same instant. An ideal instantaneous reading is the zero-duration special case. Missed events, bit errors, or unknown sensor lag require an additional observation model and are not erased by the proof.
\end{enumerate}
All experiments are local and isolated; no global failed configuration needs to be produced in the operating network. Faithful copying and clamp invariance are substantive closure assumptions. They cannot be certified against arbitrary hidden mechanisms by this finite protocol. The mission starts after acquisition and the selected preparation. If the live network keeps operating during acquisition, its exposure during that entire interval must be included in a separate safety analysis; this result does not prevent failures that precede the decision. A requirement to keep it safe throughout acquisition therefore needs a priced shutdown or holding procedure, or a proved exposure bound for that phase.

\paragraph{Control errors for deployment.}
For clarity, most formulas use an exactly specified operating field. If its stable realized value $u_i^{\rm true}$ is known only within $\xi_u$ of the reported value, and rational approximation of its node factor has log error at most $\tau_u$, replace throughout
\begin{equation}
 \begin{aligned}
 E=(N+2M)\eta&\quad\hbox{by}\quad E_u=(N+2M)\eta+N(\xi_u+\tau_u),\\
 a=(1+2d)\eta&\quad\hbox{by}\quad a_u=(1+2d)\eta+\xi_u+\tau_u.
 \end{aligned}
 \label{eq:field-error}
\end{equation}
The kinetics are measured at the same reproducible physical operating setting. This handles calibrated fixed error, not an unmodeled time-varying field.

\section{A local energy measurement without stationary samples}
\label{sec:balance}
\begin{lemma}[Finite holds, conditional concentration]
Fix a clamped two-state chain with equilibrium success probability $q$ and total rate $a\ge r_{\rm cal}$. Observe its state $Y_k$ after successive holds of length $s$. Whatever its initial state,
\[
 \mathbb E[Y_k\mid\mathcal F_{k-1}]=q+(Y_{k-1}-q)e^{-as}.
\]
If $e^{-r_{\rm cal}s}\le\beta$, then
\begin{equation}
 \Pr\left(\left|\frac1n\sum_{k=1}^nY_k-q\right|>\beta+t\right)\le2e^{-2nt^2}.
 \label{eq:martingale}
\end{equation}
\end{lemma}
\begin{proof}
The two-state transition semigroup has the displayed form. Each conditional expectation differs from $q$ by at most $\beta$. The differences $Y_k-\mathbb E[Y_k\mid\mathcal F_{k-1}]$ are martingale differences with conditional range of length one. Conditional Hoeffding's exponential bound, iterated over $k$, gives $2e^{-2nt^2}$ for their average. Adding the bounded conditional-mean bias proves the result. No independence or equilibrium preparation is used.
\end{proof}

Let $D=N+2M$ be the number of contexts consisting of the all-zero neighbor assignment at each node and each single-one neighbor assignment. Take $0<\eta\le1/4$, let
\[
 J_0=\left\lceil\log_2\max\{2,2B_0\}\right\rceil+1,
 \quad n_0=\left\lceil512\log\frac{4DJ_0}{\delta_E}\right\rceil,
 \quad n_f=\left\lceil512\eta^{-2}\log\frac{4D}{\delta_E}\right\rceil.
\]
Here and below logarithms and integer ceilings are evaluated with conservative finite rounding.

For one unknown $\ell\in[-B_0,B_0]$, maintain an interval $[L,U]$. At its midpoint $m$, command $z=-m$. Hold for $s_0=r_{\rm cal}^{-1}\log32$ between each of $n_0$ readings. The lemma gives a $1/16$ error bound for their mean $\bar q$, except on an allocated error event. If $\bar q<3/8$, replace $U$ by $m$; if $\bar q>5/8$, replace $L$ by $m$. Otherwise stop. Also stop once $U-L\le1$, using its midpoint.

On the concentration event, a lower decision gives $q<7/16$. Therefore $\ell-m+e_z<\logit(7/16)<-1/4$, where $|e_z|\le\eta/16\le1/64$, so $\ell<m$. The upper decision is analogous. Thus all retained intervals contain $\ell$. An interior stop gives
\[
 q\in[5/16,11/16],\qquad |\ell-m+e_z|\le\log(11/5)<1.
\]
A width stop gives $|\ell-m+e_z|\le1/2+1/64<1$. At most $J_0$ stages occur. The final realized logit lies in $[-1,1]$.

At that final command, take $n_f$ more readings with hold length
\[
 s_f=r_{\rm cal}^{-1}\log(32/\eta).
\]
Equation~\eqref{eq:martingale} gives $|\bar q-q|\le\eta/16$. Because $q\in[\sig(-1),\sig(1)]$ and $\eta\le1/4$, both $q$ and $\bar q$ lie in $[1/4,3/4]$. The derivative of $\logit$ there is at most $16/3<6$. Hence the estimate
\[
 \widetilde\ell=\logit(\bar q)-z
\]
has error at most $3\eta/8+\eta/16=7\eta/16$. Computing the logarithm to additive error $\eta/16$ gives
\begin{equation}
 |\widetilde\ell-\ell|\le\eta/2.
 \label{eq:local-energy}
\end{equation}
A union bound over all contexts and adaptive bisection stages gives joint success at least $1-\delta_E$. Conditional on the prior data, the held chain obeys the same concentration lemma, so adaptivity does not change the bound.

The number of bit readings is at most $D(J_0n_0+n_f)$, and total stochastic exposure is at most
\begin{equation}
 T_E\le\frac D{r_{\rm cal}}
 \left[J_0n_0\log32+n_f\log(32/\eta)\right].
 \label{eq:energy-cost}
\end{equation}
Conservative timer rounding adds its actual finite excess. Bisection's midpoint displacements form a geometric series; within one context their total is $O(B_0)$. Moving between contexts and returning to zero still costs $O(DB_0/v)$ in total slew time. All field calibration and context-setting operations are added to~\eqref{eq:energy-cost}. The algorithm does not use observations during field motion or actuator settling.

\section{Rational model reconstruction}
\label{sec:reconstruct}
Set $\widetilde h_i=\widetilde\ell_i(0)$. For each ordered interaction, form
\[
 \widetilde J_{ij}^{(i)}=\widetilde\ell_i(e_j)-\widetilde\ell_i(0),
 \qquad \widetilde J_{ij}=\tfrac12(\widetilde J_{ij}^{(i)}+\widetilde J_{ji}^{(j)}).
\]
Equation~\eqref{eq:local-energy} implies $|\widetilde h_i-h_i|\le\eta/2$ and $|\widetilde J_{ij}-J_{ij}|\le\eta$. Choose positive rational node and pair factors $w_i,w_{ij}$ satisfying
\[
 |\log w_i-\widetilde h_i|\le\eta/2,
 \qquad |\log w_{ij}-\widetilde J_{ij}|\le\eta.
\]
The learned factorization is
\[
 \widehat W(x)=\prod_i w_i^{x_i}\prod_{\{i,j\}\in\calE}w_{ij}^{x_ix_j}.
\]
Positive dyadic approximation of the exponentials gives these factors using $O(B_0+\log(1/\eta))$ bits per entry and polynomial bit work. There is no inverse probability factor, but numerical precision is polynomial in the energy magnitude $B_0$, not its binary description length alone.

With $\widehat h_i=\log w_i$ and $\widehat J_{ij}=\log w_{ij}$, the joint confidence event implies
\begin{equation}
 \begin{aligned}
 |\widehat h_i-h_i|&\le\eta, & |\widehat J_{ij}-J_{ij}|&\le2\eta,\\
 |\log\widehat W(x)-\log W(x)|&\le E=(N+2M)\eta,\\
 |\widehat\ell_i(s)-\ell_i(s)|&\le a=(1+2d)\eta.
 \end{aligned}
 \label{eq:errors}
\end{equation}
These statements cover every configuration even though the calibration queried only $D$ local contexts. This is where the known pairwise form is essential. A general unknown local table would require all its contexts.

\section{Operating kinetics from fast escape experiments}
\label{sec:kinetics}
Biased kinetics need not equal operating kinetics. For each $u,i,s$ to be learned, compute the learned operating logit $\widehat\ell_i^u(s)$. Assume the certified conditional logit error $a_u\le1/2$, refining energy measurements if necessary. Select state zero if the learned logit is nonnegative and state one otherwise. Its true exit probability factor is at least $\sig(-a_u)>1/3$. Therefore its operating exit rate $\lambda$ lies in $[r_*/3,r^*]$.

At the operating field, pin that state, release, and record its first-exit dwell. Repeat, charging every pin, release, settling interval, and read. The strong Markov and faithful-reset assumptions make the post-release dwells independent $\operatorname{Exp}(\lambda)$ variables. They are actively generated dwell experiments, not free equilibrium samples.

Let
\[
 C=\sum_{u\in\calU}\sum_{i=1}^N2^{\deg(i)},\quad
 k=\left\lceil256\rho^{-2}\log\frac{8C}{\delta_Q}\right\rceil,\qquad
 T_{\rm cap}=\frac3{r_*}\log\frac{4Ck}{\delta_Q},\quad 0<\rho\le1/4.
\]
Abort this attempted certificate if a dwell exceeds its cap. From quantized dwell sum $\widetilde S$, set $\widehat\lambda=k/\widetilde S$.

\begin{lemma}[Finite-resolution kinetic measurement]
With probability at least $1-\delta_Q$, all requested measurements finish within their caps and obey $|\log(\widehat\lambda/\lambda)|\le\rho$. Total dwell exposure is at most $CkT_{\rm cap}$; expected uncapped exposure is at most $3Ck/r_*$. Pinning and settling add $k\sum_{u,i,s}C_{\rm pin}(i,s,u)$.
\end{lemma}
\begin{proof}
For $S=\sum_{j=1}^kT_j$, $\lambda S$ is gamma with shape $k$ and unit rate. Chernoff's moment-generating-function calculation gives
\[
 \Pr(|\lambda S/k-1|>t)\le2e^{-kt^2/4},\qquad0<t<1.
\]
Take $t=\rho/8$ and union bound over $C$ experiments. Timestamp error gives $|\widetilde S-S|\le2k\Delta$, so
\[
 |\lambda\widetilde S/k-1|\le\rho/8+2r^*\Delta\le5\rho/32<\rho/4.
\]
For $\rho\le1/4$ the logarithm of the reciprocal differs from zero by less than $\rho$. The union probability that any of $Ck$ dwells exceeds $T_{\rm cap}$ is at most $Ck\exp(-r_*T_{\rm cap}/3)\le\delta_Q/4$. The stated choices leave more than enough confidence budget. Exposure bounds follow from $\lambda\ge r_*/3$.
\end{proof}

If the measured direction is $0\to1$, set
\[
 \widehat Q_i^u(0\to1\mid s)=\widehat\lambda,
 \qquad \widehat Q_i^u(1\to0\mid s)=\widehat\lambda e^{-\widehat\ell_i^u(s)}.
\]
If it is $1\to0$, reverse the roles and multiply by $e^{\widehat\ell_i^u(s)}$. The exponential ratios are products of the learned rational factor entries. Thus the rate table is rational. Local detailed balance makes the learned generator reversible with respect to the learned Gibbs law. The measured direction has log error at most $\rho$ and the other at most
\begin{equation}
 b_u=\rho+a_u.
 \label{eq:rate-error}
\end{equation}
In particular the rare operating direction is inferred from a measured fast escape and an experimentally acquired energy difference, with its uncertainty retained.

This step enumerates $K\sum_i2^{\deg(i)}$ contexts. That cost is explicit and can be exponential in $d$. If the physically justified operating model instead has a context-independent clock $r_i^u$, one fast-direction measurement per $i,u$ suffices to estimate it; in that reduced version conservatively use rate log error $\rho+2a_u$. Across different operating fields their clocks are still measured separately unless invariance is independently part of the input model. Stationary-risk comparison alone does not require kinetic acquisition.

\section{From data to verifiable premises}
\label{sec:checks}
The following are finite interval checks on the learned model and its confidence envelope. A failed check returns ``premises unresolved''; it does not justify using an uncertified algorithm.

For each operating field enumerate the local contexts. Compute the learned Dobrushin influences $\widehat A_{ij}$. Since $\sig$ is $1/4$-Lipschitz and each logit error is at most $a_u$,
\begin{equation}
 A_{ij}^{\rm true}\le\widehat A_{ij}+a_u/2\quad(j\in\partial i).
 \label{eq:influence}
\end{equation}
Use $\widehat A_{ij}+a_u/2$ as an outward upper bound. Check that each row sums to at most $1-\alpha_u$ for a positive certified $\alpha_u$. The learned and true models then both satisfy that certificate. All pinnings inherit it.

For a visible coordinate, let $\widehat L_i^u$ and $\widehat U_i^u$ be its learned extreme logits. The true conditional failure probabilities lie in
\[
 q_i^- =\sig(\widehat L_i^u-a_u),\qquad q_i^+ =\sig(\widehat U_i^u+a_u).
\]
With outward rational rounding, check $q_i^-\ge(q_i^+)^c$ for a fixed chosen integer $c$. No minimization over unknown global states is hidden here; the local context enumeration is charged.

The rate bands give explicit positive clock bounds:
\[
 \begin{aligned}
 \underline r_u&=\min_{i,s}\left(\widehat Q_i^u(0\to1\mid s)+\widehat Q_i^u(1\to0\mid s)\right)e^{-b_u},\\
 \overline r_{i,u}&=\max_s\left(\widehat Q_i^u(0\to1\mid s)+\widehat Q_i^u(1\to0\mid s)\right)e^{b_u}.
 \end{aligned}
\]
For the true generator, Dirichlet-form comparison with rate-one heat bath gives gap at least $\gamma_u=\alpha_u\underline r_u$; the same bounds conservatively cover the learned generator. A valid uniformization rate is $R_u=2\sum_i\overline r_{i,u}$, and $g_u=\gamma_u/R_u$ is a supplied gap certificate. State-dependent clocks cause no problem: subtracting the constant heat-bath backbone of rate $\underline r_u$ leaves nonnegative reversible flip rates.

The accepted static solver gives a confidence interval for the learned $\widehat p_u$. Widen it by $e^{\pm2E_u}$ and check that its upper endpoint is at most $1/2$. This certifies the true safe mass required for physical preparation. If necessary refine the model/solver, with an explicit confidence allocation.

\paragraph{When success is guaranteed.}
The construction is always sound on its confidence event. Guaranteed production of a certificate requires margin from the premise boundaries. For example, suppose each true influence row is at most $1-2\alpha_0$, its true surprisal contrast is at most $c_0<c$, and $p_u\le1/2-\mu$. For sufficiently small $a_u$, $d a_u\le\alpha_0$ makes the influence check pass. Writing $f(z)=-\log\sig(z)$, direct differentiation gives $|(\log f)'|\le1$. The outer logit envelope is at most $2a_u$ beyond each true extremum; consequently its contrast is at most $c_0e^{4a_u}$ and passes whenever $4a_u<\log(c/c_0)$. Sufficiently accurate relative static bounds pass the safe-mass check. Dependence on $\alpha_0^{-1}$, $\mu^{-1}$, and the contrast margin is retained. No finite-data test proves an exact boundary equality.

\section{Finite physical preparation with a relative law}
\label{sec:prep}
Fix an operating field whose checks passed. Let $B_u$ be an outward bound on every true full-conditional logit magnitude. Moving between two configurations by at most $N$ bit flips gives
\begin{equation}
 \pi_u(x)\ge p_{\min,u}:=\exp\{-N(B_u+\log2)\}\quad\hbox{for every }x.
 \label{eq:pimin}
\end{equation}
Indeed the ratio between any two unnormalized weights is at most $e^{NB_u}$, and there are $2^N$ states. This is a conservative, explicit energy bound, not the catastrophic probability.

For a reversible finite chain with gap $\gamma_u$, the standard eigen-expansion of the mean-zero density gives
\[
 \TV(P_t(x,\cdot),\pi_u)\le\tfrac12\pi_u(x)^{-1/2}e^{-\gamma_ut}
 \le\tfrac12p_{\min,u}^{-1/2}e^{-\gamma_ut}.
\]
For completeness, the initial chi-square norm is $1/\pi_u(x)-1$; reversibility contracts its square by $e^{-2\gamma_ut}$; Cauchy--Schwarz bounds total variation by half its square root.

Given $0<\zeta\le1/4$, evolve the real system for
\begin{equation}
 T_u=\gamma_u^{-1}\left[\tfrac32N(B_u+\log2)+\log(1/\zeta)\right]
 \label{eq:prep-time}
\end{equation}
from its current arbitrary state. Round time upward to the instrument's resolution. Then for every starting distribution and every state $y$,
\begin{equation}
 (1-\zeta)\pi_u(y)\le P_{T_u}(x,y)\le(1+\zeta)\pi_u(y).
 \label{eq:pointwise}
\end{equation}
The total variation bound is at most $\zeta p_{\min,u}$ and hence bounds each singleton error by at most $\zeta\pi_u(y)$.

Inspect the configuration after a known deterministic inspection duration $C_{\rm inspect}$ and accept if its completion-time state lies in $A$. The system continues to evolve throughout inspection; the extra deterministic duration preserves~\eqref{eq:pointwise}, since a Markov semigroup preserves a pointwise sandwich around its invariant law. Otherwise continue evolving for another $T_u$ and inspect again. No reset is required. Since $\pi_u(A)\ge1/2$, every attempt succeeds with conditional probability at least $(1-\zeta)/2\ge3/8$. After
\[
 L_{\rm prep}=\left\lceil\log(1/\delta_{\rm prep})/\log(8/5)\right\rceil
\]
attempts, abort if none was safe. The abort probability is at most $\delta_{\rm prep}$ and elapsed time is at most $L_{\rm prep}(T_u+C_{\rm inspect})$. This is a law for the randomized accepted state; it is not a claim about the risk conditional on each realized state after its reading.

\begin{lemma}[Prepared-law sandwich]
Conditional on preparation succeeding, its accepted state law $\nu_u^{\rm prep}$ satisfies, for every safe state,
\begin{equation}
 e^{-\beta}\pi_u(x\mid A)\le\nu_u^{\rm prep}(x)\le e^\beta\pi_u(x\mid A),
 \qquad \beta=\log\frac{1+\zeta}{1-\zeta}.
 \label{eq:start-sandwich}
\end{equation}
Every future path-event probability under these two starts consequently differs by at most $e^{\pm\beta}$, including the first-failure CDF for every $t\ge0$.
\end{lemma}
\begin{proof}
For each attempt, conditional on any earlier history, the starting law is arbitrary but~\eqref{eq:pointwise} still holds. Conditioning its endpoint on $A$ divides by a number between $(1-\zeta)\pi_u(A)$ and $(1+\zeta)\pi_u(A)$, proving the displayed ratio. The law at the first successful capped attempt is a mixture of such conditional laws; the ratio bounds survive mixtures. A future path event has probability $\sum_x\nu(x)k(x)$ for a nonnegative conditional probability $k(x)$, so the same pointwise sandwich transfers to it.
\end{proof}

The appearance of $p_{\min,u}$ does not cause inverse-$p_{\min,u}$ elapsed time: only its logarithm enters~\eqref{eq:prep-time}. Slow physical relaxation is paid through $\gamma_u^{-1}$. This construction requires a repairable Markov system. A preparation procedure that terminates forever upon its first catastrophe has a different accepted law; the present lemma must not be applied to that killed procedure.

\subsection{An optional bound on preparation's own failure exposure}
If the operating dynamics are attractive in the coordinate order and $B$ is increasing, starting at all-zero gives a useful extra bound. Constant-clock ferromagnetic heat bath is one sufficient case. Arbitrary context-dependent clocks are not automatically attractive; that additional condition must be proved or checked separately.

Couple the all-zero-started process below a stationary process. For any elapsed preparation horizon $T$, its chance of ever entering $B$ is at most the stationary process's chance. Here $T$ includes inspections and begins after the priced all-zero initialization and operating-field setup; any exposure during those earlier operations must be charged separately. Stationarity and a jump-count bound imply
\[
 \Pr_\pi(\exists t\le T:X_t\in B)
 \le p_u+T\sum_{x\in A}\pi_u(x)Q^u(x,B).
\]
Reversibility makes the boundary flux equal to $\sum_{x\in B}\pi_u(x)Q^u(x,A)\le p_u\Lambda_u$, where $\Lambda_u=\max_x\sum_{y\ne x}Q^u(x,y)\le\sum_i\overline r_{i,u}$. Thus
\begin{equation}
 \Pr_0(\hbox{a catastrophe during preparation by }T)\le p_u(1+T\Lambda_u).
 \label{eq:prep-exposure}
\end{equation}
Use $T=L_{\rm prep}T_u$ and the certified upper bound on $p_u$ to obtain a data-derived exposure bound. Initializing the whole system at zero has its own priced clamp/setup cost. This estimate can make preparation exposure small; it is not a promise that no catastrophe can occur.

\section{Central acquisition-to-certificate theorem}
\begin{theorem}[Costed active acquisition, preparation, and prevention]
\label{thm:central}
Under Section~\ref{sec:access}, the preceding finite protocol, followed by the accepted solver, has the following property. Except with the allocated sum of experimental and computational error probabilities, it either reports unresolved premises, or returns, for each checked operating field $u$:
\begin{enumerate}
\item a positive rational Gibbs model, a rational reversible local generator, and a simultaneous confidence envelope covering their true factor values and every true transition rate;
\item a relative interval for stationary disconnection probability;
\item one finite evaluable certificate $[L_u(t),U_u(t)]$ valid for the actual prepared-start first-failure CDF simultaneously at every positive rational deadline; and
\item a finite physical preparation procedure whose abort probability is at most $\delta_{\rm prep}$ and whose elapsed cost is at most $L_{\rm prep}(T_u+C_{\rm inspect})$, with $T_u$ given by~\eqref{eq:prep-time}.
\end{enumerate}
If the accepted solver's learned-model CDF interval has ratio at most $1+\eps_0$, the returned interval may be taken as
\begin{equation}
 L_u(t)=e^{-(6E_u+b_u+\beta)}\widehat L_u(t),\qquad
 U_u(t)=\min\{1,e^{6E_u+b_u+\beta}\widehat U_u(t)\},
 \label{eq:final-band}
\end{equation}
with ratio at most $(1+\eps_0)e^{2(6E_u+b_u+\beta)}$. Arbitrarily prescribed relative width $1+\eps$ is achieved by choosing positive $\eta,\rho,\zeta$, field precision, factor rounding, and $\eps_0$ so that
\[
 \log(1+\eps_0)+2(6E_u+b_u+\beta)\le\log(1+\eps)
\]
for all candidates, subject to the premise margins.

For any deadline $t_*>0$, possibly selected after seeing the simultaneous certificates but before observing the selected physical preparation's accepted state, choose a candidate minimizing $U_u(t_*)$. Its true prepared-start failure probability is at most $(1+\eps)$ times the smallest true prepared-start probability in the checked menu. The same assertion applies to stationary risk using its static intervals.

Elapsed acquisition and deployment cost is the sum of~\eqref{eq:energy-cost}, the actual certified slew cost for total field travel $V_{\rm total}$ across all phases, all stated cell/context/environment/field/read/pin costs, at most $CkT_{\rm cap}$ dwell exposure, and the selected whole-system preparation and inspection cost. Energy calibration alone has commanded field travel $O(DB_0)$, with realized travel additionally bounded by the readback intervals; movements between operating settings and deployment are also included in $V_{\rm total}$. Actuator work and apparatus-construction resources are separate ledger entries, with no implicit conversion into elapsed time. Computational work includes explicit local context enumeration $C=K\sum_i2^{\deg(i)}$, rational interval arithmetic, and the accepted solver's polynomial cost at fixed contrast $c$. No factor $1/\pi_u(B)$ or mean catastrophe lifetime is introduced by acquisition or preparation. If actuator or mixing costs themselves have such a dependence, that dependence is fully retained.
\end{theorem}
\begin{proof}
The energy and kinetic lemmas give the common confidence event and finite exposure. Equations~\eqref{eq:errors}--\eqref{eq:rate-error} give the factor and common-support rate envelopes. Section~\ref{sec:checks} verifies the accepted solver's hypotheses and the physical gap/safe-mass bounds. The inherited comparison~\eqref{eq:inherited} gives multiplicative CDF error $e^{6E_u+b_u}$ between the true and learned equilibrium-safe laws. The prepared-law lemma contributes $e^\beta$ uniformly over all future path events. Widening the solver's simultaneous interval proves~\eqref{eq:final-band}. A union bound over the finite menu covers all computational runs and permits deadline selection after certificate construction and before observing the selected prepared state. Only the selected candidate must be physically prepared; the other menu laws specify hypothetical applications of the same preparation protocol.

For selection, let $u_*$ minimize the actual prepared-start CDF and let $\widehat u$ minimize the reported upper endpoint. Then
\[
 F_{\widehat u}^{\rm prep}(t_*)\le U_{\widehat u}(t_*)\le U_{u_*}(t_*)
 \le(1+\eps)L_{u_*}(t_*)\le(1+\eps)F_{u_*}^{\rm prep}(t_*).
\]
The hypothetical preparation law for an unselected candidate is defined by the stated procedure conditional on successful preparation; no physical run of that candidate is required for the comparison. All operations and costs have already been enumerated. Positive error budgets of order $\eps/(N+M+d)$ for energy and $\eps$ for rate and initialization suffice; the sample formulas are polynomial in their reciprocals. Strict premise margins guarantee that the finite checks pass at sufficiently small accuracy, as shown above.
\end{proof}

\paragraph{One explicit accuracy allocation.}
For $0<\eps\le1$, let $s=\log(1+\eps)$. In the ideal exactly specified operating-field case, choose
\[
 \eta\le\frac{s}{48(N+2M)+8(1+2d)},\quad
 \rho\le s/16,\quad \zeta\le s/64,\quad \log(1+\eps_0)\le s/4.
\]
Then $6E+a\le s/8$, $\rho\le s/16$, and $\beta\le3\zeta\le3s/64$, so the total logarithmic width is at most $s/4+2(15s/64)=23s/32<s$. Additional margin is available for field/rounding errors. Also enforce the premise margins and $a_u\le1/2$. The constants are deliberately sufficient, not optimized.

\section{Why insufficient field amplitude forces rare exposure}
\label{sec:lower}
The previous advantage is purchased by control access. A one-coordinate experiment already provides an obstruction.

\begin{theorem}[Adaptive amplitude--exposure lower bound]
Consider a two-state chain with known attempt clock $r$ and unknown logit $\theta$. An adaptive experiment may use the entire continuous jump trajectory, arbitrary predictable fields $|z_t|\le H$, arbitrary parameter-independent forced state resets, and any independent randomization. Resets reveal no parameter-dependent work, force, or other side-channel measurement; their only output is the commanded state. Total time during which the chain is exposed is at most $T$. Start from a parameter-independent state. If an estimate distinguishes $\theta_0=-B_0$ from $\theta_1=-B_0+4\eta$, with error at most $\delta<1/2$ under both and $0<\eta\le1/8$, then
\begin{equation}
 T\ge\frac{(1-2\delta)\log((1-\delta)/\delta)}{64r\eta^2}
 e^{(B_0-H)_+}.
 \label{eq:lower}
\end{equation}
The same bound applies to estimating the logit to additive error $\eta$ on both instances.
\end{theorem}
\begin{proof}
At field $z$ let $q_j=\sig(\theta_j+z)$ and $\Delta_\theta=4\eta\le1/2$. The two exit rates are $rq_j$ from zero and $r(1-q_j)$ from one. For positive rates $a,b$, the Poisson divergence rate satisfies
\[
 a\log(a/b)-a+b\le(a-b)^2/b.
\]
Since $q_1-q_0=(e^{\Delta_\theta}-1)q_0(1-q_1)$ and $e^{\Delta_\theta}-1\le2\Delta_\theta$, either starting state has divergence rate at most $4r\Delta_\theta^2q_0\le64r\eta^2q_0$. Also $q_0\le\min\{1,e^{-B_0+H}\}=e^{-(B_0-H)_+}$.

The controlled-path likelihood chain rule integrates this bound over exposure time. The policy's actions have the same conditional kernels under both parameters and add no divergence; forced resets likewise add none. Thus the transcript laws obey
\[
 \KL(P_0\Vert P_1)\le64rT\eta^2 e^{-(B_0-H)_+}.
\]
A successful binary test has a decision event whose probabilities are at least $1-\delta$ and at most $\delta$. Data processing to that event lower-bounds the KL divergence by $(1-2\delta)\log((1-\delta)/\delta)$. Additive-$\eta$ estimates on parameters separated by $4\eta$ give such a test by thresholding at their midpoint. Combining the two inequalities proves the claim.
\end{proof}

For the following slew refinement, let $T$ denote total elapsed wall time, including field motion and intervals when the system is pinned or not observed. If the field starts at zero and has finite slew $|\dot z|\le v$, a sharper available-information bound is
\begin{equation}
 \KL(P_0\Vert P_1)\le64r\eta^2\int_0^T
 \min\{1,e^{-B_0+\min(H,vt)}\}\,dt.
 \label{eq:slew}
\end{equation}
For $H\le B_0$, the integral is exactly
\[
 e^{-B_0}\left[\frac{e^{v\min(T,H/v)}-1}{v}+e^H(T-H/v)_+\right].
\]
Hence sufficient amplitude with prohibitively slow control is not free access. The lower bound gives an exponential field-deficit dependence; the positive construction gives polynomial exposure once the compensating scale is accessible. We do not claim its polynomial constants or its multi-context total cost are minimax optimal.

\section{Executed synthetic interface diagnostic}
\label{sec:diagnostic}
The accompanying \texttt{diagnostic.py} implements the stated access restrictions on a synthetic four-bit system. The visible network has four parallel edges, so all four must fail to disconnect it. Its interaction graph is a cycle, with $h_i=-28$ and $J_{ij}=8$. Clocks explicitly depend on both neighbor context and field. The learner receives finite-time local bit observations and measured, quantized prepared first-exit dwell times; hidden parameters are inspected only by the final audit. There is no input stream of global catastrophe occurrences.

The two supplied operating settings are $u=0$ and $u_i=-3$. The diagnostic uses $\eta=0.04$, $\rho=0.1$, $\delta_E=\delta_Q=0.02$, and a fixed random seed. All local factor/rate audit errors satisfy their proposed bounds. The conservative influence checks and fixed contrast $c=3$ pass. Exact enumeration of the 16 global states is used only to expose the reference static risk and to exercise the confidence-envelope arithmetic; the accepted general rare-CDF solver is not reimplemented or benchmarked here.

\begin{center}
\begin{tabular}{lcc}
\toprule
 & Baseline & Protective setting\\
\midrule
Reference stationary risk & $1.80485\times10^{-35}$ & $1.10894\times10^{-40}$\\
Widened lower endpoint & $6.92892\times10^{-36}$ & $4.25728\times10^{-41}$\\
Widened upper endpoint & $4.72619\times10^{-35}$ & $2.90387\times10^{-40}$\\
Maximum audited rate log error & $0.004541$ & $0.004090$\\
Certified gap diagnostic & $0.4748$ & $0.4914$\\
\bottomrule
\end{tabular}
\end{center}
The widened intervals are disjoint and select the protective setting. These are relatively coarse intervals: $E=0.48$ and $b=0.3$, so the model-to-CDF distortion bound is $6E+b=3.18$. This numerical run is not an example of a narrow CDF certificate. The theorem's smaller accuracy allocation is available at its stated larger acquisition cost.

The serial specified-component ledger contains 30,226,360 local bit readings, 7,747,360 forced preparations, 32 kinetic contexts, 44 context settings, and 100 field commands. The maximum realized bias magnitude is $28.126$ to the displayed precision. The measured energy-hold exposure is $4.01922\times10^8$ time units and the operating-dwell exposure is $7.53586\times10^6$. Including the declared per-pin, per-context, per-read, event-detection, and available-slew costs gives $4.09497107\times10^8$ time units. If the rates are interpreted in inverse seconds, that is approximately 12.98 years on one serial experimental channel. The absence of inverse-global-rarity cost is a mathematical scaling result, not a demonstrated practical acquisition speed.

The selected field's planned physical preparation has a capped elapsed bound of about 3,988 additional time units, including specified inspections. It is calculated, not physically executed. Cell construction, instrument calibration, full menu-environment setup, deployment slew, and actuator work are not assigned numerical prices in this simulator; they remain additional explicit terms in the theorem's ledger. The checks use ordinary floating arithmetic, not the outward interval arithmetic required for an instrument-certified implementation. Exact reproducible data, the seed, and all component costs are in \texttt{diagnostic\_results.json}.

\section{Primary prior work and the defensible delta}
Designed perturbations of spin networks are established prior work. Jiang, Sivak, and Thomson \cite{jiang} actively choose fields using Fisher information to improve inverse-Ising inference; their experiments include large sample reductions. That directly anticipates the qualitative idea of breaking strongly correlated or uninformative regimes by fields. Our elementary bisection is not presented as the first active spin-network method.

Dutt, Lokhov, Vuffray, and Misra \cite{dutt} study parameter/structure learning from Glauber-type dynamics, including multi-start single-step data and active initialization; their theoretical and empirical work shows large gains away from equilibrium. Bresler, Gamarnik, and Shah \cite{bresler} establish efficient learning from observed Glauber trajectories. These are the appropriate dynamical-learning comparators. Their learning objectives should not be relabeled as rare-CDF certificates, and our stronger priced clamp/release interface should not be relabeled as their observational access model.

Scarlett and Cevher \cite{scarlett} prove lower bounds for an active graphical-model selection model that chooses which nodes to observe. Their selection access differs from changing the physical conditional law. The lower bound in Section~\ref{sec:lower} is therefore proved directly for controlled continuous-time fields and resets, instead of importing that result across access models.

Bennett's free-energy estimator \cite{bennett} and Torrie--Valleau umbrella sampling \cite{torrie} establish the broad importance of overlap and biasing for rare thermodynamic weights. The present construction acquires an unknown local model before invoking a rare-event computation; it is not a new invention of importance sampling.

The defensible contribution of this document is the complete, explicit composition: finite physical local measurements; measured operating kinetics without cross-field clock transport; uncertainty strong enough for relative rare-event prediction; finite relative safe initialization; finite-menu intervention choice; and an amplitude/exposure boundary. Historical novelty, optimality, and superiority over all prior methods remain unestablished. The retrieved primary sources suffice to prevent broad novelty claims, not to prove a literature-wide absence.

\section{Limitations that change the conclusion}
\begin{enumerate}
\item The known complete interaction graph and pairwise, Markov, reversible intervention model are essential. Unknown common causes, missing variables, aging, state-dependent temperature, or nonlocal kinetic channels can invalidate extrapolation to catastrophe states.
\item Calibrated fields and state pins are resources. Their construction, readback, work, slew, settling, and rare-state holding costs are charged. Polynomial statistical cost does not prove that a natural system admits a cheap actuator.
\item If only node observation can be chosen, or $H\ll B_0$, the positive access model is unavailable; Section~\ref{sec:lower} quantifies one resulting obstruction.
\item Arbitrary local clocks require all $2^{\deg(i)}$ contexts at every dynamic operating candidate. A sparse parameter graph does not by itself make unknown kinetic tables sparse.
\item Finite physical initialization applies to a repairable process. Irreversible failure during setup, censoring of setup paths, or a required all-working operating start needs a different theorem. Optional attractiveness gives a setup exposure bound, not a replacement initial law.
\item Prevention is a certified comparison of a finite supplied menu. The menu's size and acquisition cost are paid. No solution to arbitrary continuous, combinatorial, or hidden-constraint intervention design is claimed.
\item The synthetic diagnostic accompanying this source checks the acquisition interface and cost accounting on a known simulator. It is not evidence for faithful physical clamping in a real network, and it does not revalidate the accepted global CDF solver.
\end{enumerate}

\begin{thebibliography}{9}
\bibitem{accepted} Accepted Astra construction, \emph{A finite first-failure curve for interacting networks}, 5 November 2026, and its execution supplement. Local immutable inputs: \texttt{next3\_root/correlated\_reliability\_report.tex} and \texttt{next3\_root/execution\_supplement.tex}. Public corpus at fixed revision \texttt{fef3640a8b1dc5475674a006de417f652f3baf3c}, \href{https://github.com/harjassand/astra-research-knowledge/tree/fef3640a8b1dc5475674a006de417f652f3baf3c/research/returns/2026-11-05-p3-correlated-reliability}{permanent public corpus link}.
\bibitem{prior} Accepted governing generic CDF proof, local \texttt{next2\_reliability/PROOF.md}; preserved as supplied.
\bibitem{jiang} J. Jiang, D. A. Sivak, and M. Thomson, \emph{Active Learning of Spin Network Models}, arXiv:1903.10474 (2019), \url{https://arxiv.org/abs/1903.10474}; author repository \url{https://authors.library.caltech.edu/records/cwa0t-yys04}.
\bibitem{dutt} A. Dutt, A. Y. Lokhov, M. Vuffray, and S. Misra, \emph{Exponential Reduction in Sample Complexity with Learning of Ising Model Dynamics}, Proceedings of ICML, PMLR 139:2914--2925 (2021), \url{https://proceedings.mlr.press/v139/dutt21a.html}.
\bibitem{bresler} G. Bresler, D. Gamarnik, and D. Shah, \emph{Learning Graphical Models from the Glauber Dynamics}, IEEE Transactions on Information Theory 64(6):4072--4080 (2018); original Allerton 2014 manuscript, \url{https://arxiv.org/abs/1410.7659}.
\bibitem{scarlett} J. Scarlett and V. Cevher, \emph{Lower Bounds on Active Learning for Graphical Model Selection}, Proceedings of AISTATS, PMLR 54:55--64 (2017), \url{https://proceedings.mlr.press/v54/scarlett17a.html}.
\bibitem{bennett} C. H. Bennett, \emph{Efficient estimation of free energy differences from Monte Carlo data}, Journal of Computational Physics 22:245--268 (1976), \url{https://doi.org/10.1016/0021-9991(76)90078-4}.
\bibitem{torrie} G. M. Torrie and J. P. Valleau, \emph{Nonphysical sampling distributions in Monte Carlo free-energy estimation: Umbrella sampling}, Journal of Computational Physics 23:187--199 (1977), \url{https://doi.org/10.1016/0021-9991(77)90121-8}.
\end{thebibliography}
\end{document}
